\documentclass[11pt]{article}
\usepackage[margin=1in]{geometry}
\usepackage{amsmath,amssymb,mathtools}
\usepackage{booktabs}
\usepackage{enumitem}
\usepackage{siunitx}
\sisetup{group-separator={,},group-minimum-digits=3}

\title{\textbf{Solutions -- Quiz 1} \vspace{ 0.4cm }\\ECON 300: Labour Economics}
\author{Instructor: Muhebullah Karimzada}


\begin{document}
\maketitle

\section*{Multiple Choice (5 points)}


\begin{enumerate}
\item \textbf{Answer: (c)}. A discouraged worker who stopped searching is counted as \textbf{not in the labour force} because they are not actively searching for work. In standard labour force concepts, active search is required to be classified as unemployed.

\item \textbf{Answer: (b)}. The employment-to-population ratio is:
\[
\frac{E}{\text{Working-age population}}.
\]
It measures the share of the working-age population that is employed, regardless of whether others are unemployed or out of the labour force.

\item \textbf{Answer: (c)}. An increase in non-labour income raises consumption at every leisure choice without changing the wage. Graphically, it shifts the budget line upward in a parallel way in the labour--leisure diagram.

\item \textbf{Answer: (c)}. At sufficiently high wages, the income effect can dominate the substitution effect, leading people to choose more leisure and fewer work hours. That generates a backward-bending labour supply curve over that range.

\item \textbf{Answer: (b)}. A more generous unemployment insurance system can lengthen search duration or reduce the urgency of accepting a job, and can also affect measured participation through search behaviour, potentially lowering measured participation if some people reduce active search or delay re-entry. (The direction can depend on program design, but among the options, this is the best fit.)
\end{enumerate}

\newpage

\subsection*{Problem 1: Labour Force Measurement and Discouragement (7.5 points)}

\textbf{Given initial conditions:}
\[
\text{Working-age population} = 900{,}000,\quad
E = 585{,}000,\quad
U = 65{,}000,\quad
N = 250{,}000.
\]

\subsubsection*{(a) Labour force, LFPR, unemployment rate (2.5 points)}

\textbf{Step 1: Labour force.}
\[
LF = E + U = 585{,}000 + 65{,}000 = 650{,}000.
\]

\textbf{Step 2: LFPR.}
\[
LFPR = \frac{LF}{\text{Working-age population}}
= \frac{650{,}000}{900{,}000}
= 0.7222\ldots \approx 72.22\%.
\]

\textbf{Step 3: Unemployment rate.}
\[
u = \frac{U}{LF}
= \frac{65{,}000}{650{,}000}
= 0.10 = 10\%.
\]

\subsubsection*{(b) Employment-to-population ratio and interpretation (2.5 points)}

\textbf{Step 1: Employment-to-population ratio.}
\[
\frac{E}{\text{Working-age population}}
= \frac{585{,}000}{900{,}000}
= 0.65 = 65\%.
\]

\textbf{Step 2: Conceptual difference from LFPR.}
The employment-to-population ratio measures the share of the working-age population that is actually employed.
The LFPR measures the share of the working-age population that is \emph{engaged with the labour market}, meaning employed or actively searching (unemployed).
Therefore, LFPR can be high even if employment is not high, as long as many people are searching (unemployed).

\subsubsection*{(c) Slowdown: job losses and discouragement (2.5 points)}

\textbf{Event 1: 15,000 employed lose jobs.} Assume they initially move to unemployment (they want work).
\[
E' = 585{,}000 - 15{,}000 = 570{,}000,
\qquad
U' = 65{,}000 + 15{,}000 = 80{,}000.
\]

\textbf{Event 2: 10,000 unemployed stop searching.} They move from unemployed to not in the labour force.
\[
U'' = 80{,}000 - 10{,}000 = 70{,}000.
\]

Now the labour force is:
\[
LF'' = E' + U'' = 570{,}000 + 70{,}000 = 640{,}000.
\]

\textbf{New unemployment rate:}
\[
u'' = \frac{U''}{LF''}
= \frac{70{,}000}{640{,}000}
= 0.109375 \approx 10.94\%.
\]

\textbf{Explanation.}
Discouragement reduces measured unemployment because individuals who stop searching are no longer counted as unemployed; they exit the labour force.
This changes both:
\begin{itemize}
\item the numerator $U$ (falls when people stop searching), and
\item the denominator $LF$ (also falls when people stop searching).
\end{itemize}
As a result, the unemployment rate may not rise as much as the true deterioration in job opportunities would suggest. This is why economists also examine participation rates and employment-to-population ratios when evaluating labour market conditions.

\subsection*{Problem 2: Labour Supply and Minimum Wage (7.5 points)}

\textbf{Given individual decision:}
\[
C = 500 + 30h,
\qquad
U = C - h^2.
\]

\textbf{Given market:}
\[
L^D = 150 - 3w,
\qquad
L^S = 30 + 2w.
\]

\subsubsection*{(a) Optimal hours of work (2.5 points)}

\textbf{Step 1: Substitute the budget constraint into utility.}
\[
U(h) = (500 + 30h) - h^2
= 500 + 30h - h^2.
\]

\textbf{Step 2: Take the derivative and set to zero (FOC).}
\[
\frac{dU}{dh} = 30 - 2h.
\]
FOC:
\[
30 - 2h = 0
\quad \Rightarrow \quad
2h = 30
\quad \Rightarrow \quad
h^* = 15.
\]

\textbf{Step 3: Second-order condition.}
\[
\frac{d^2U}{dh^2} = -2 < 0,
\]
so the solution is a maximum.

\textbf{Answer:} The worker supplies \textbf{15 hours per week}.

\subsubsection*{(b) Competitive equilibrium wage and employment (2.5 points)}

\textbf{Step 1: Set $L^D = L^S$.}
\[
150 - 3w = 30 + 2w.
\]

\textbf{Step 2: Solve for $w$.}
\[
150 - 30 = 2w + 3w
\quad \Rightarrow \quad
120 = 5w
\quad \Rightarrow \quad
w^* = 24.
\]

\textbf{Step 3: Compute equilibrium employment.}
\[
L^* = 150 - 3(24) = 150 - 72 = 78.
\]
(Check: $30 + 2(24) = 30 + 48 = 78$.)

\textbf{Answer:} Competitive equilibrium is \textbf{$w^* = 24$} and \textbf{$L^* = 78$}.

\subsubsection*{(c) Minimum wage $w = 25$ (2.5 points)}

\textbf{Step 1: Is it binding?} Compare to equilibrium:
\[
w_{\min} = 25 > w^* = 24,
\]
so the minimum wage is \textbf{binding}.

\textbf{Step 2: Labour demanded at $w=25$.}
\[
L^D(25) = 150 - 3(25) = 150 - 75 = 75.
\]

\textbf{Step 3: Labour supplied at $w=25$.}
\[
L^S(25) = 30 + 2(25) = 30 + 50 = 80.
\]

\textbf{Step 4: Unemployment (excess supply).}
\[
U = L^S - L^D = 80 - 75 = 5.
\]

\textbf{Answer:} The minimum wage is \textbf{binding} and generates \textbf{5} units of unemployment (in the same units as $L$, e.g., thousands if $L$ is in thousands).

\end{document}
